%% Własne makra i ustawienia dokumentu.

% Przykładowe makro skrótu
\newcommand{\np}{np.\@\xspace}

% Polskie nazwy dla \cref (cleveref nie ma wbudowanego języka polskiego)
\crefname{section}{rozdz.}{rozdz.}
\Crefname{section}{Rozdział}{Rozdziały}
\crefname{subsection}{podrozdz.}{podrozdz.}
\Crefname{subsection}{Podrozdział}{Podrozdziały}
\crefname{figure}{rys.}{rys.}
\Crefname{figure}{Rysunek}{Rysunki}
\crefname{table}{tab.}{tab.}
\Crefname{table}{Tabela}{Tabele}
\crefname{equation}{wzór}{wzory}
\Crefname{equation}{Wzór}{Wzory}
\crefformat{equation}{wzór~(#2#1#3)}
\crefname{appendix}{zał.}{zał.}
\Crefname{appendix}{Załącznik}{Załączniki}

% Kolumna tabularx z wyrównaniem do lewej
\newcolumntype{L}{>{\raggedright\arraybackslash}X}

% Łagodniejsze łamanie wierszy z długimi słowami i adresami URL
\setlength{\emergencystretch}{2.5em}
\AtBeginBibliography{\sloppy}
